\chapter{Data and Preprocessing}
\label{ch:data}

\section{Requirements for Evaluating an Agentic IR System}
\label{sec:data-requirements}

A dataset is chosen because it can answer the questions a report has to ask.
This system's central claim, that a validation agent can send useful feedback
to a planning agent, is meaningful only if it can be falsified with numbers.
Four requirements follow, stated before any dataset is named.

\textbf{Queries with relevance judgements.} Recall@$k$, nDCG@10 and MRR are
defined over a mapping from queries to relevant documents. Without one those
metrics are undefined, and the comparison against baseline IR systems that
Chapter~\ref{ch:implementation} must present cannot be computed.

\textbf{Genuinely compositional needs.} On single-fact look-ups the correct
decomposition is the identity, every plan has one node, and the agentic
apparatus collapses onto the hybrid baseline it is supposed to beat. Measuring
a planner needs questions no single document answers.

\textbf{Entity and structure annotations.} Sentence-level supporting facts make
retrieved evidence checkable at the granularity the system emits; gold
relational triples make a graph traversal itself checkable.

\textbf{A corpus that indexes and runs on consumer hardware.} Every component
runs on a single 8\,GB laptop GPU with no paid API. A web-scale corpus is not
expensive under that constraint; it is unusable.

BEIR~\cite{thakur2021beir} satisfies the first and fourth requirements across
eighteen tasks, but most of its tasks express single-hop needs with no
reasoning-structure annotations. The two datasets adopted below satisfy all
four, and the first of them is itself one of BEIR's tasks.

\section{On the Datasets Listed in the Assignment Brief}
\label{sec:table2}

Table~2 of the brief lists The Pile, Common Crawl, WikiText, OpenWebText and
LAION-5B. All five are pretraining corpora: unlabelled text, or weakly
labelled image and text pairs, at a scale suitable for fitting a model. None
ships a query set, a relevance judgement or a gold answer, so the first
requirement fails outright. This is not a deficiency in those corpora but a
statement of what they are for. They are, in fact, the right resources for a
different layer of this system. The 8B model that does the reasoning here owes
its grasp of English, of world entities and of the shape of a JSON object to
pretraining on exactly such corpora. The brief's list describes the substrate
the agent brings to the task; it does not describe the benchmark retrieval is
scored against. The brief's own instruction to check for updated datasets is
the licence to select evaluation data on evaluation criteria.

\section{HotpotQA}
\label{sec:hotpotqa}

HotpotQA~\cite{yang2018hotpotqa} is a multi-hop question answering dataset over
English Wikipedia in which each question requires facts from two distinct
paragraphs. This work uses the \texttt{distractor} setting, where each
question carries its two gold paragraphs and eight distractors, and the
\texttt{validation} split of 7{,}405 questions.

Its distinguishing property is the supporting-facts annotation: every question
carries the \texttt{(title, sent\_id)} pairs a human annotator used. This is
retrieval ground truth at sentence granularity, and it fixes the granularity of
the rest of the system. Evidence is carried as sentences, so that what the
Verifier reasons over and what the metrics score are the same objects.

Reading the contexts of all 7{,}405 questions yields 73{,}700 paragraphs, which
reduce to 66{,}581 distinct passages after deduplication; 7{,}119 discarded
paragraphs were byte-identical repeats and no title carries two body texts. The
corpus holds 269{,}602 sentences, a mean of 4.049 per passage. Of 18{,}005 gold
supporting facts, 18{,}004 resolve; the exception cites sentence 902 of a
five-sentence paragraph, an upstream annotation error reported rather than
dropped, because the same counter would catch a real identifier drift.

One measured property governs the sampling design. HotpotQA labels questions
easy, medium or hard, but every one of the 7{,}405 validation questions is
\texttt{hard}; the easier levels exist only in training. The informative
partition is reasoning type: 5{,}918 bridge questions, where one entity must be
resolved before the second hop can be posed, and 1{,}487 comparison questions,
where two entities are retrieved and an attribute compared.

\section{2WikiMultihopQA}
\label{sec:twowiki}

2WikiMultihopQA~\cite{ho2020constructing} was built to expose the reasoning
steps a question requires, not only its endpoints. Its validation split holds
12{,}576 questions whose 125{,}760 context paragraphs reduce to 54{,}957
distinct passages and 202{,}759 sentences. It ships 30{,}687 sentence-level
supporting facts, all of which resolve, and, in addition, 31{,}120 gold
Wikidata evidence triples covering every question. Those triples convert the
KG Navigator from a component this report describes into one it measures.

They also introduce the most serious leakage risk in the project. A graph built
from those triples would let the KG Navigator read the answer key. The graph is
therefore built from passage text only, the triples are loaded exclusively by
scoring code, and a test asserts that the graph builder never opens a file
whose name contains \texttt{evidence} or \texttt{triple}.

All four reasoning types are populated: compositional (5{,}236), comparison
(3{,}040), bridge-comparison (2{,}751) and inference (1{,}549). The third is
the hardest case in either dataset, requiring a bridge resolution and a
comparison over its results. The split is also textually messier: because the
same article appears in many contexts with inconsistent detokenisation, 1{,}720
titles ship two or more variants, producing 2{,}617 title collisions, and since
\texttt{sent\_id} indexes a sentence list the variant kept decides whether a
gold fact lands on the right sentence.

\section{Preprocessing Pipeline}
\label{sec:preprocessing}

The pipeline is deliberately thin. It applies only transformations that cannot
be undone downstream and declines every one a consumer can perform for itself.

\paragraph{Segmentation and identifiers.} Both datasets ship each paragraph as
a sentence list, and that list is adopted verbatim. No sentence splitter is
called anywhere, because the list is the identifier space of the ground truth:
re-splitting would renumber every gold fact while the metrics went on producing
plausible numbers. Each passage is keyed by \texttt{source:quote(title)}, so a
retrieved passage, a graph node and a gold fact name the same object the same
way.

\paragraph{Deduplication.} HotpotQA collapses cleanly. 2Wiki leaves 2{,}617
cases where one title carries two texts; the rule keeps the variant with the
most sentences, the widest identifier space, and its cost is measured: keeping
the first occurrence would leave 70 gold identifiers out of range and 211 on
the wrong sentence, whereas keeping the widest leaves none out of range and 150,
or 0.489\%, on a slightly different detokenisation of the intended sentence.

\paragraph{What is deliberately not done.} The only text transformation is the
collapse of whitespace runs. There is no lowercasing, stopword removal,
stemming or punctuation stripping at the corpus level, because those are
properties of a scorer's analyser: BM25 stems inside its own analyser, the
dense encoder needs original case to distinguish proper nouns, and the NLI
model needs well-formed sentences. One convention is shared by all scorers:
every one sees the same title-prefixed unit, \texttt{"Title. body"}, because
fusing a channel that indexes titles with one that does not produces rankings
that cannot be compared.

\paragraph{Entity graph and indexes.} Nodes are normalised passage titles
whose aliases are compiled into a trie for longest-match linking. A directed
edge $A \to B$ is created when an alias of $B$ occurs in the text of $A$, and
the edge stores the sentence asserting it, which is what makes a traversal
result citable at the same granularity as a retrieved passage. The HotpotQA
graph retains 64{,}733 nodes and 123{,}759 edges; the 2Wiki graph 52{,}959
nodes and 63{,}073 edges. A document-frequency cap removed 446 and 172 hub
nodes, titles such as \texttt{one} whose occurrences are not relational
assertions. Each corpus is also indexed twice, by BM25 with Snowball stemming
and by an exact FAISS index over \texttt{bge-small-en-v1.5} embeddings; the
two occupy 129.1\,MB and 103.6\,MB, small enough to hold in memory, which is
why exact search is used and no recall is lost at the index level.

\section{Evaluation Sampling}
\label{sec:sampling}

Evaluation runs on a frozen 250-question sample per dataset, and the reason is
arithmetic. At 52 tokens per second a full agentic question was estimated
before the runs at 22 to 25 seconds; the medians actually measured came out
higher, and the estimate is recorded as the estimate it was. On it, one agentic
configuration over the full 7{,}405-question split is about 45 hours, and the
grid is nine configurations over two datasets. At 250 questions the grid was
feasible on this hardware; at 7{,}405 it was not.

The sample is drawn once with seed 42, materialised to disk, and never
re-sampled; its SHA-256 is recorded in every run's metadata so a silently
changed sample is detectable afterwards. Stratification is on question type,
not difficulty, and that is a measured decision: with every HotpotQA
validation question labelled hard, stratifying on level would give a single
stratum, an unstratified sample presented as a stratified one. Stratifying on
type preserves the bridge and comparison ratio exactly at 200 and 50
questions, and 2Wiki's four types at 104, 60, 55 and 31.
Table~\ref{tab:dataset-sample} gives the composition.

Two hundred and fifty questions is enough to separate the systems compared
here. Treated as an unpaired proportion near 0.5, the 95\% half-width is about
6.2 points. The comparisons that matter are paired: every configuration
answers the identical 250 questions, so the bootstrap is taken over
per-question differences, and the correlation between systems sharing a
corpus, an index and a question set makes those intervals materially tighter.
Each dataset also gets a disjoint 50-question calibration slice, verified to
share no question identifier with the 250, reserved for setting the Verifier's
threshold. Table~\ref{tab:dataset-corpus} summarises the two corpora as built.

% Chapter 2 -- data and preprocessing
% GENERATED by src/agentic_ir/eval/tables.py -- do not edit by hand.
% Regenerate with: python -m agentic_ir.cli tables
% Requires \usepackage{booktabs} and \usepackage{graphicx};
% report/main.tex already loads both.
% Source runs:
%   hotpotqa/agentic_full: agentic_full_hotpotqa_20260901T214116Z (5 questions)

\begin{table}[htbp]
  \centering
  \small
  \setlength{\tabcolsep}{4pt}
  \caption{Corpus statistics after preprocessing, read from \texttt{\{dataset\}\_corpus\_stats.json}. Passages are de-duplicated by content under a canonical title-derived id; sentence splits are the datasets' own and are persisted verbatim, because gold supporting facts index into them by position and re-splitting anywhere downstream would renumber the ground truth while every metric still looked plausible. \texttt{--} marks a statistic the build did not record.}
  \label{tab:dataset-corpus}
  \resizebox{\ifdim\width>\textwidth\textwidth\else\width\fi}{!}{%
  \begin{tabular}{lrr}
    \toprule
    Statistic & \texttt{hotpotqa} & \texttt{twowiki} \\
    \midrule
    Questions in split & 7\,405 & 12\,576 \\
    Paragraphs read & 73\,700 & 125\,760 \\
    Passages after de-duplication & 66\,581 & 54\,957 \\
    Identical paragraphs dropped & 7\,119 & 68\,186 \\
    Title collisions resolved & 0 & 2\,617 \\
    Sentences & 269\,602 & 202\,759 \\
    Mean sentences per passage & 4.05 & 3.69 \\
    Mean passage length (chars) & 531.9 & 427.2 \\
    Mean sentence length (chars) & 130.6 & 115.1 \\
    Questions with gold triples & 0 & 12\,576 \\
    Gold supporting facts & 18\,005 & 30\,687 \\
    Supporting facts unresolvable & 1 & 0 \\
    Evaluation sample & 250 & 250 \\
    \bottomrule
  \end{tabular}}
\end{table}

\begin{table}[htbp]
  \centering
  \small
  \setlength{\tabcolsep}{4pt}
  \caption{Composition of the frozen evaluation slice against the full split. The slice is stratified on question type and sampled once, with seed 42, by \texttt{scripts/sample\_eval\_set.py}; its SHA-256 is recorded in every run's \texttt{meta.json} and it is never re-sampled at run time. HotpotQA's validation split is entirely \texttt{level=hard}, so difficulty gives one stratum and a silently unstratified sample; question type is the informative split there.}
  \label{tab:dataset-sample}
  \resizebox{\ifdim\width>\textwidth\textwidth\else\width\fi}{!}{%
  \begin{tabular}{lrrrr}
    \toprule
    Question type & Split $n$ & Split share & Sample $n$ & Sample share \\
    \midrule
    \multicolumn{5}{@{}l}{\textit{\texttt{hotpotqa}}} \\
    bridge & 5\,918 & 0.799 & 200 & 0.800 \\
    comparison & 1\,487 & 0.201 & 50 & 0.200 \\
    \textit{total} & 7\,405 & 1.000 & 250 & 1.000 \\
    \addlinespace
    \multicolumn{5}{@{}l}{\textit{\texttt{twowiki}}} \\
    bridge\_comparison & 2\,751 & 0.219 & 55 & 0.220 \\
    comparison & 3\,040 & 0.242 & 60 & 0.240 \\
    compositional & 5\,236 & 0.416 & 104 & 0.416 \\
    inference & 1\,549 & 0.123 & 31 & 0.124 \\
    \textit{total} & 12\,576 & 1.000 & 250 & 1.000 \\
    \bottomrule
  \end{tabular}}
\end{table}

